\section{Exact two-mark clearance flux on one physical flight}
\label{sec:dc-two-mark}
The preceding clearance--roof analysis treats a selected near miss as a
\emph{mark}, without adding a collision.  Two different unvisited disks
can be marked on the same actual flight.  Here we quantify their
physical multiplicity exactly, rather than assuming independence or
multiplying by the number of collision words.  The result is a
second-order \emph{collision-space} clearance theorem; it is not an
$m^2$-normalized pointwise roof estimate.

Throughout let $R\in[R_-,R_+]$, $R_+=47/100$, and keep
$\delta_*=3/200$, $0<\delta_1,\delta_2\le\delta_*$, the triangular
lattice $\Lambda$, the actual flight time $\tau\le\tau_+$, and the
one-copy physical mark maps $X_R,Z_R$ of
Theorem~\ref{thm:cf-flux}.  Define the strictly positive \emph{physical
separation} and the finite comparison set
\begin{align}
 c_{\rm pair}
   &=\sqrt{1-[2(R_++\delta_*)]^2}
      =\frac{\sqrt{591}}{100}>0,\label{eq:dc-constants}\\
 \mathcal W_*
   &=\{w\in\Lambda\setminus\{0\}:|w|\le\tau_++2(R_++\delta_*)\},
   \notag\\
 N_{\rm pair}&=|\mathcal W_*|.\notag
\end{align}
The constant $N_{\rm pair}$ is finite by the \emph{same} finite-horizon
assumption already used in Lemma~\ref{lem:cf-foot}; it is not the
number of admissible words, the number of pairing sheets, or an
unproved distortion constant.

\begin{lemma}[Distinct same-flight clearances are uniformly transverse]
\label{lem:dc-transverse}
Let $z_1\ne z_2$ be two unvisited disk centers whose clearance values
$u_i\in(0,\delta_i)$ are realized by the \emph{same} regular physical
flight with unit direction $v_\vartheta$.  Let
$b_i=(z_i-q)\cdot v_\vartheta^\perp$ be their signed impact parameters.
Then both perpendicular feet lie in the open physical flight,
$z_2-z_1\in\mathcal W_*$, and
\begin{equation}\label{eq:dc-transverse}
 |(z_2-z_1)\cdot v_\vartheta|
   =\sqrt{|z_2-z_1|^2-(b_2-b_1)^2}
   \ge c_{\rm pair}.
\end{equation}
In particular, for the first signed impact parameter held fixed,
$\partial_\vartheta b_2=-(z_2-z_1)\cdot v_\vartheta$ does not vanish
on a simultaneous physical two-mark state.
\end{lemma}
\begin{proof}
Lemma~\ref{lem:cf-foot} puts each foot in the open physical
flight.  Thus its along-flight coordinate lies in an interval of
length at most $\tau_+$.  Since $|b_i|=R+u_i$,
\[
 |b_2-b_1|\le 2R+\delta_1+\delta_2
       \le2(R_++\delta_*)=97/100<1.
\]
The lattice spacing $|z_2-z_1|\ge1$ and orthogonal decomposition
along $v_\vartheta,v_\vartheta^\perp$ imply the equality and lower
bound in \eqref{eq:dc-transverse}.  The same decomposition gives
$|z_2-z_1|\le\tau_++2(R_++\delta_*)$, proving membership in
$\mathcal W_*$.  The impact derivative follows from
$d v_\vartheta^\perp/d\vartheta=-v_\vartheta$ with $b_1$ fixed.
Neither a grazing endpoint nor an unvisited near miss is treated as
an extra collision.
\end{proof}

For a first center translated to $0$, fix $w\in\mathcal W_*$ and
$\sigma_1,\sigma_2\in\{\pm1\}$.  Write
\begin{equation}\label{eq:dc-delta}
 d_{\sigma_1\sigma_2}(u_1,u_2)
   =\sigma_2(R+u_2)-\sigma_1(R+u_1).
\end{equation}
For all $0<u_i<\delta_i$,
$|d_{\sigma_1\sigma_2}|<97/100<|w|$.  Thus the equation
\begin{equation}\label{eq:dc-angle-root}
 w\cdot v_\vartheta^\perp
     =d_{\sigma_1\sigma_2}(u_1,u_2)
\end{equation}
has exactly two roots modulo $2\pi$, denoted
$\vartheta_{w,\sigma_1,\sigma_2,+}(u_1,u_2)$ and
$\vartheta_{w,\sigma_1,\sigma_2,-}(u_1,u_2)$, distinguished by
the sign of $w\cdot v_\vartheta$.  They are smooth on each
parameter rectangle.  Their one-dimensional change-of-variables
Jacobian is
\begin{equation}\label{eq:dc-pair-jacobian}
 \left|\frac{\partial\vartheta}{\partial u_2}\right|
  =\frac1{\sqrt{|w|^2-d_{\sigma_1\sigma_2}(u_1,u_2)^2}},
\end{equation}
with $u_1$ fixed; the sign of $\sigma_2$ has unit magnitude.

\begin{theorem}[An exact physical two-mark coarea identity]
\label{thm:dc-exact}
For every nonnegative Borel weight $a(x,z_1,z_2)$ of a genuine
outgoing collision state and two \emph{ordered}, distinct unvisited
centers (expressed in the initial incident disk's lattice frame),
set at a root of \eqref{eq:dc-angle-root}
\[
 x=X_R(\vartheta,\sigma_1,u_1),\quad
 z_1=Z_R(\vartheta,\sigma_1,u_1),\quad z_2=z_1+w,
 \quad
 \mathfrak A=\one_{\{z_2\in\mathcal Z_{\delta_2}(x)\}}.
\]
Then the \emph{full, one-copy} original collision measure obeys
\begin{align}\label{eq:dc-exact}
 &\int\sum_{z_1\in\mathcal Z_{\delta_1}(x)}
       \sum_{\substack{z_2\in\mathcal Z_{\delta_2}(x)\\ z_2\ne z_1}}
         a(x,z_1,z_2)\,d\nu(x)\\
 &\qquad=\frac1{4\pi R}
 \sum_{w\in\mathcal W_*}
 \sum_{\sigma_1,\sigma_2,\eta\in\{\pm1\}}
 \int_0^{\delta_1}\int_0^{\delta_2}
  \left.
   \frac{\mathfrak A\,a(x,z_1,z_1+w)}
        {\sqrt{|w|^2-d_{\sigma_1\sigma_2}(u_1,u_2)^2}}
  \right|_{\vartheta=\vartheta_{w,\sigma_1,\sigma_2,\eta}}
   du_2\,du_1.\notag
\end{align}
The mask $\mathfrak A$ retains \emph{both actual feet}, the same
physical collision flight, and the prohibition on unvisited-disk
reflections.  It must not be removed in an equality.  Any original
return, displacement, collision, roof, section, or radius test may
be included in $a$ as a Borel selector; \eqref{eq:dc-exact}
never modifies that test or the physical clock.
\end{theorem}
\begin{proof}
Apply the proved one-mark flux identity
\eqref{eq:cf-flux} to the first mark $z_1$.  This gives the factor
$(4\pi R)^{-1}$, the sign $\sigma_1$, the $u_1$ integral and the
original angular integral.  For each $(\vartheta,u_1)$, translating
$z_1$ to the origin associates any distinct second center with
exactly one $w=z_2-z_1\in\mathcal W_*$ by
Lemma~\ref{lem:dc-transverse}.  The second center has a definite
sign $\sigma_2$ of its impact and a unique $u_2\in(0,\delta_2)$;
its physical membership is precisely $\mathfrak A=1$.
For fixed $w,\sigma_1,u_1$, the function
$\vartheta\mapsto b_2=b_1+w\cdot v_\vartheta^\perp$ is a sine wave
around the fixed $b_1$.  On each of the two signed impact bands
$\sigma_2(R,R+\delta_2)$, equation \eqref{eq:dc-angle-root} has
exactly two roots, and its derivative has the nonzero magnitude
in \eqref{eq:dc-pair-jacobian}.  One-dimensional substitution on
the two monotone branches therefore gives the sum over $\eta$ and
the displayed denominator.  The total sum over $w$ is finite;
Tonelli and monotone approximation first for bounded weights then
for arbitrary nonnegative Borel $a$ justify all substitutions.
Only regular charts of the genuine flight are used; the inherited
one-mark construction assigns a fixed null-set convention on
singular lines.  The formula counts an ordered pair once because
its first mark is unique in the first flux identity and its
relative $w$ and sign/angle root are unique thereafter.  No
physical word count, symbolic renewal process, conditional
independence, or fictitious collision is introduced.
\end{proof}

\begin{theorem}[Uniform quadratic bound for two distinct unvisited marks]
\label{thm:dc-factorial}
For the original invariant collision measure and every pair of
allowed clearances,
\begin{equation}\label{eq:dc-factorial}
 \int \sum_{z_1\in\mathcal Z_{\delta_1}(x)}
   \sum_{\substack{z_2\in\mathcal Z_{\delta_2}(x)\\z_2\ne z_1}}
                 1\,d\nu(x)
 \le\frac{2N_{\rm pair}}{\pi R c_{\rm pair}}\,\delta_1\delta_2
 \le\frac{2N_{\rm pair}}{\pi R_-c_{\rm pair}}\,\delta_1\delta_2.
\end{equation}
This bound includes every physical line/flight and both signs; it
does not require independence between the two marks.
\end{theorem}
\begin{proof}
Apply Theorem~\ref{thm:dc-exact} with $a\equiv1$.  Drop only the
positive \emph{mask in the upper bound}.  The denominator is at
least $c_{\rm pair}$ by the lattice-spacing calculation, and there
are exactly $N_{\rm pair}$ candidate differences and $2^3=8$
sign/root choices.  Thus the right side is at most
$8N_{\rm pair}\delta_1\delta_2/(4\pi Rc_{\rm pair})$, which is
\eqref{eq:dc-factorial}.  The denominator lower bound is pointwise
on the entire parameter rectangle, not only almost everywhere
on the physically admitted pairs.
\end{proof}

\begin{corollary}[A sharp first-order unique-clearance law]
\label{cor:dc-unique}
Let $N_\delta=|\mathcal Z_\delta|$ be the \emph{original} one-flight
unvisited-near-miss count and write
$C_{\rm pair}=N_{\rm pair}/(\pi R_-c_{\rm pair})$.  Then
\begin{align}
 0&\le\frac{\delta}{R}-\nu(N_\delta\ge1)
                    \le C_{\rm pair}\delta^2,
                                      \label{eq:dc-union}\\
 0&\le\nu(N_\delta\ge2)
                    \le C_{\rm pair}\delta^2,\label{eq:dc-multiple}\\
 \nu(N_\delta=1)&\ge\frac{\delta}{R}
                       -2C_{\rm pair}\delta^2.
                                      \label{eq:dc-exact-one}
\end{align}
In particular $\nu(N_\delta\ge1)=\delta/R+O(\delta^2)$,
uniformly in the original radius interval; the $\delta/R$ term is
exactly the one-copy physical flux, not a proposed Poisson law.
\end{corollary}
\begin{proof}
For every nonnegative integer $N$,
$0\le N-\one_{\{N\ge1\}}=(N-1)_+\le N(N-1)/2$ and
$\one_{\{N\ge2\}}\le N(N-1)/2$.  Also
$N\one_{\{N\ge2\}}\le N(N-1)$.  Apply these deterministic
inequalities to $N_\delta$, use the exact physical mean
\eqref{eq:cf-exact-mean}, and substitute
Theorem~\ref{thm:dc-factorial} with
$\delta_1=\delta_2=\delta$.  No limit in collision count is
needed for these one-flight statements.
\end{proof}

\begin{corollary}[Original exact-label records: uniform weak overcount payment]
\label{cor:dc-schedule}
Fix $m,R,n,k$, the original physical $m$-flight record, and a
nonnegative Borel roof test $0\le\psi\le1$.  Let
$0<\delta_j\le\delta_*$ for $0\le j<m$.  With the original
$\Pi=\Pi_{n,k,m,R}$ and unchanged normalized source $\nu_R^*$,
\begin{align}\label{eq:dc-original-record}
 0&\le\int\one_{\Pi}\psi(L)
  \sum_{j=0}^{m-1}\bigl[
       N_{\delta_j}(T_R^jx)
     -\one_{\{N_{\delta_j}(T_R^jx)\ge1\}}
                   \bigr]d\nu_R^*(x)\\
 &\le\frac{C_{\rm pair}}{c_*}
             \sum_{j=0}^{m-1}\delta_j^2.\notag
\end{align}
For the original depth schedule
$\delta_j=2\varepsilon r^{d'_j}$,
$0<2\varepsilon\le\delta_*$, $r=q_0^{1/4}$,
this is at most
\begin{equation}\label{eq:dc-scheduled}
        \frac{8C_{\rm pair}\varepsilon^2}
                        {c_*(1-r^2)},
\end{equation}
uniformly in $m,R,n,k$ and the test $\psi$.  Thus the positive
roof measure of within-flight \emph{duplicate marking} has
uniform total variation of order $\varepsilon^2$, even with the
full original exact labels.  The corresponding density need not
have controlled $L^\infty$ height: multiplying this mass estimate
by $m^2$ at fixed outer $B$ is \emph{not} justified.
\end{corollary}
\begin{proof}
The integrand in \eqref{eq:dc-original-record} is nonnegative.
Since $\Pi\subset Y_R^*$ and $\nu_R^*=c_*^{-1}\one_{Y_R^*}\nu$,
its integral is bounded by $c_*^{-1}$ times the sum of unconditioned
collision integrals with $\psi\le1$.  By invariance of the original
collision measure, each integral equals the same one-flight
quantity at $T_R^jx$ and is bounded by
$C_{\rm pair}\delta_j^2$ from
Corollary~\ref{cor:dc-unique}.  The depth list has at most two
indices at each nonnegative depth, so
$\sum_j\delta_j^2\le
2(2\varepsilon)^2\sum_{d\ge0}r^{2d}
=8\varepsilon^2/(1-r^2)$.  All source and time tests were
retained until applying positive domination; no independent
probability factor and no exchange of the original limits occurs.
\end{proof}

The exact second-order formula \eqref{eq:dc-exact} isolates the
nontransverse configurations which a two-mark pointwise analysis
would otherwise need to pay.  The separation constant
$c_{\rm pair}>0$ rules out that particular simultaneous same-flight
angle degeneracy on the original radius range.  This leaves
\emph{single} unvisited-disk clearance, marks on different
physical flights, grazing/section/pairing boundary terms, and the
nongrazing capacity shortage.  In particular the full
$\mathcal H_M(\mathsf R_B^{\rm low})$ limit and the original raw
pointwise local limit theorem remain open; the new $O(\varepsilon^2)$
total variation gain is never substituted for the required
$m^2L^\infty$ estimate.
